% !TEX program = xelatex
\documentclass[a4paper]{article}
\usepackage{amsthm}
\usepackage{amssymb}
\usepackage{bm}
\usepackage{mathtools}
\usepackage[x11names]{xcolor}
\usepackage{xparse}
\usepackage{fontspec}
\usepackage{unicode-math}
\setromanfont{DovesType-Regular.otf}
\setsansfont{Andika}
\setmathfont{Asana Math}[Scale=1]

% \usepackage{pstricks}
\usepackage{varwidth}
\usepackage{siunitx}
\usepackage{graphicx}
\usepackage[margin=1cm]{geometry}
\usepackage[most]{tcolorbox}
\usepackage{pgfplots}
\pgfplotsset{compat=newest}
\tcbuselibrary{skins,xparse,poster,breakable}
% \usetikzlibrary{fadings}
\usetikzlibrary{calc, plotmarks, shapes, shapes.geometric, positioning, angles, intersections, quotes, through, patterns, turtle, arrows.meta}
\usetikzlibrary{decorations.markings,backgrounds}
% \usepackage{etoolbox}
% \usepackage{tkz-euclide}
% \usepackage{xlop}
% \newcommand\hole[2]{#1}  % for use with xlop
\pagenumbering{gobble}
%%%%%%%%%%%%%%%%%%%%%%%%%%%%%%%%%%%%%%%%%%%%%%%%%%%%%%%%%
\newcommand\markangle[9]{% origin X Y radius radiusmark mark colour opacity
%  % fill red circle offset-from-centre
  \begin{scope}
    \path[clip] (#1) -- (#2) -- (#3);
    \fill[color=#7,fill opacity=#8,draw=black,name path global=pcircle]  % global declaration required otherwise pcircle is not seen by the `named intersections=' lines below.
    (#1) circle (#4);
  \end{scope}
  % middle calculation
  \path[name path=line one] (#1) -- (#2);
  \path[name path=line two] (#1) -- (#3);
  \path[%
  name intersections={of=line one and pcircle, by={inter one}},
  name intersections={of=line two and pcircle, by={inter two}}
  ] (inter one) -- (inter two) coordinate[pos=#9] (place);
  % put mark
  \node at ($(#1)!#5!(place)$) {\scriptsize{#6}};
}

\newcommand\tcircle[6]{% centre coord (x,y), radius, points, radpoint, colour, edge
  \coordinate (O) at (#1,#2); % centre of the circle
  \def\radius{#3}          % radius of the circle
  \def\npts{#4}            % number of the points
  \def\radpt{#5}           % radius of the points
  \colorlet{ptcolour}{#6}  % colour of the points
  % \draw (O) circle (\radius);
  \foreach \numpoint in {1,...,\npts}{
    \fill[ptcolour] (O) ++ (360/\npts*\numpoint:\radius) coordinate (C\numpoint) circle(\radpt);
  }
}

% \newcommand{\condSoln}[2]{\ifcsdef{r@#1}{#2}{}}

% \newcommand\fadingtext[3][]{%
%    \begin{tikzfadingfrompicture}[name=fading letter]
%      \node[text=transparent!0,inner xsep=0pt,outer xsep=0pt,#1] {#3};
%    \end{tikzfadingfrompicture}%
%    \begin{tikzpicture}[baseline=(textnode.base)]
%      \node[inner sep=0pt,outer sep=0pt,#1](textnode){\phantom{#3}};
%      \shade[path fading=fading letter,#2,fit fading=false]
%      (textnode.south west) rectangle (textnode.north east);%
%    \end{tikzpicture}%
% }

\definecolor{JISpurple}{RGB}{89,72,122}
\definecolor{JISivory}{RGB}{241,234,221}
\definecolor{JIStaupe}{RGB}{183,156,154}
\definecolor{PaleGreen}{RGB}{240,255,240} % 'Honeydew'

\AddToHook{shipout/background}{%
    \put (0in,-\paperheight){\includegraphics[width=\paperwidth,height=\paperheight]{images/R10V5-R.png}}%
}

\newcommand\numberthis{\addtocounter{equation}{1}\tag{\theequation}}

\newtcolorbox{MyOuterBox}{%
  enhanced,
  breakable,
  frame style=JISpurple,
  colback=JISivory,
  colframe=JISpurple,
  title={\includegraphics[width=0.9cm,height=0.9cm]{images/JIS Final Logo FA-02.png}\raisebox{3mm}{\Large{Maths Challenge}\hspace{26em} \Large{\bfseries\sffamily 12}}},
}

\newtcolorbox{MyInnerBox}[2][]{enhanced,%empty,
coltitle=JISpurple,colback=white,
breakable,
fonttitle=\bfseries\sffamily,
attach boxed title to top left={yshift=-1.5mm},
boxed title style={empty, size=small, top=1mm, bottom=0pt},
varwidth boxed title=0.5\linewidth,
frame code={
  \path (title.east|-frame.north) coordinate (aux);
\path[draw=JISpurple, line width=0.5mm, rounded corners,fill=white]
(frame.west) |- ([xshift=-2.5mm]title.north east) to[out=0, in=180] ([xshift=7.5mm]aux)-|(frame.east)|-(frame.south)-|cycle;
},
title={#2},#1}

\newtcolorbox{MyInnerSplitBox}[2][]{enhanced,%empty,
bicolor,sidebyside,sidebyside align=top seam,
righthand width=8.5cm,colbacklower=white,
sidebyside gap=5mm,
breakable,
coltitle=JISpurple,colback=white,
fonttitle=\bfseries\sffamily,
attach boxed title to top left={yshift=-1.5mm},
boxed title style={empty, size=small, top=1mm, bottom=0pt},
varwidth boxed title=0.5\linewidth,
frame code={
  \path (title.east|-frame.north) coordinate (aux);
\path[draw=JISpurple, line width=0.5mm, rounded corners,fill=white]
(frame.west) |- ([xshift=-2.5mm]title.north east) to[out=0, in=180] ([xshift=7.5mm]aux)-|(frame.east)|-(frame.south)-|cycle;
},
title={#2},#1}


\newtcolorbox{MySolutionBox}[1][]{%
  enhanced,
  breakable,
  frame style=JISpurple,
  colback=PaleGreen, colframe=green,
  title={\Large Solution},
  drop fuzzy shadow,
  halign=left,
  #1
}

%%%%%%%%%%%%%%%%%%%%%%%%%%%%%%%%%%%%%%%%%%%%%%%%%%
\newtoggle{SOLUTION}
%%% Uncomment the appropriate line below to show solutions %%%
% \toggletrue{SOLUTION}
\togglefalse{SOLUTION}
%%%%%%%%%%%%%%%%%%%%%%%%%%%%%%%%%%%%%%%%%%%%%%%%%


%%%%%%%%%%%%%%%%%%%%%%%%%%%%%%%%%%%%%%%%%%%%%%%%%%
%%%%%%            DOCUMENT BEGINS           %%%%%%
%%%%%%%%%%%%%%%%%%%%%%%%%%%%%%%%%%%%%%%%%%%%%%%%%%
\begin{document}


  \begin{MyOuterBox}
    \iftoggle{SOLUTION}{Here are the full, or partial solutions.
    }{
      Welcome to this week's Maths Challenge!\\
      It's fine to have a go at both questions!\\
      Solutions must be explained in detail, responses with just the answer may be ignored.\\
      Drop your solution in the box in the staffroom by Tuesday.
    }
       \begin{MyInnerBox}{Year 8 and below}
     You have \SI{10}{\litre} of nard oil in a \SI{10}{\litre} alabaster jar, which you need to divide into two equal volumes. You have an empty \SI{7}{\litre} jar and an empty \SI{3}{\litre} jar. Using only the jars, measure out two jars of \SI{5}{\litre} each.
     (This puzzle first appeared in a book in Japan in 1631.)\par
      \iftoggle{SOLUTION}{%conditional output begin
      \begin{MySolutionBox}
        Let's set up our three jars with the starting conditions at Step \(0\):\par
        \begin{tabular}{c|c|c|c|l}
            Step & \SI{10}{\litre} jar & \SI{7}{\litre} jar & \SI{3}{\litre} jar\\\hline
            \(0\) & \(10\) & \(0\) & \(0\)\\
            \(1\) & \(3\) & \(7\) & \(0\) & Fill the \SI{7}{\litre} jar from the \SI{10}{\litre} jar\\
            \(2\) & \(3\) & \(4\) & \(3\) & Fill the \SI{3}{\litre} jar from the \SI{7}{\litre} jar\\
            \(3\) & \(6\) & \(4\) & \(0\) & Pour the \SI{3}{\litre} jar into the \SI{10}{\litre} jar\\
            \(4\) & \(6\) & \(1\) & \(3\) & \(\vdots\)\\
            \(5\) & \(9\) & \(1\) & \(0\)\\
            \(6\) & \(9\) & \(0\) & \(1\)\\
            \(7\) & \(2\) & \(7\) & \(1\)\\
            \(8\) & \(2\) & \(5\) & \(3\)\\
            \(9\) & \(5\) & \(5\) & \(0\)\\
          \end{tabular}\par\vspace{2mm}
        Now we have divided the \SI{10}{\litre}, \SI{5}{\litre} in the \SI{10}{\litre} jar, and \SI{5}{\litre} in the \SI{7}{\litre} jar.
      \end{MySolutionBox}
    }{}%conditional output end
    \end{MyInnerBox}


    \vspace{0.4cm}
          \begin{MyInnerBox}{Year 9 and above}
        Without calculating the answer, is
        \[
          15.99 \div 3.99
        \]
        more than \(4\), less than \(4\) or exactly \(4\)?
      \iftoggle{SOLUTION}{%conditional output begin
      \begin{MySolutionBox}
        If we think of this quotient as being
        \begin{align*}
          \frac{15.99}{3.99} & = \frac{4 + 4 + 4 + 3.99}{3.99}\\
          &= \frac{4}{3.99} + \frac{4}{3.99} + \frac{4}{3.99} + \frac{3.99}{3.99}\\
          &= slightly > 4 + slightly > 4 + slightly > 4 + 1
        \end{align*}
        So the answer must be a bit larger than \(4\).\\
      Or, we might start with \(\frac{16}{4}\) and note that to get \(\frac{15.99}{3.99}\) the numerator and the denominator have each been decreased by the same amount, \(0.01\). But proportionally the denominator has been decreased by more than the numerator (because it is smaller than the numerator) so the quotient must \emph{increase}.
      \end{MySolutionBox}
    }{}%conditional output end
    \end{MyInnerBox}


  \end{MyOuterBox}

%%%%%%%%%%%%%%%%%%%%%%%%%%%%%%%%%%%%%%%%%%%%%%%%%%
\end{document}



